\section*{致谢}

我们感谢 MILC 合作组分享用于本研究的格点组态。
LQCD 计算使用 Chroma 软件套件完成~\cite{Edwards:2004sx}。
本研究使用了国家能源研究科学计算中心（NERSC）的资源——该中心是美国能源部科学办公室用户设施，由美国能源部科学办公室通过合同 DE-AC02-05CH11231 经 ERCAP 计划支持；
还使用了 USQCD 合作组的计算设施，这些设施由美国能源部科学办公室资助；
并部分得到密歇根州立大学通过 Institute for Cyber-Enabled Research（iCER）提供的计算资源支持。
WG 的工作部分受到美国能源部科学办公室资助 DE-SC0024053 ``High Energy Physics Computing Traineeship for Lattice Gauge Theory" 以及美国国家科学基金会资助 PHY~2209424 的支持。
FY 受美国能源部科学办公室核物理办公室合同 DE-SC0012704 支持，并在 Scientific Discovery through Advanced Computing（SciDAC）项目 ``Fundamental Nuclear Physics at the Exascale and Beyond" 框架内获得支持。
HL 部分受到美国国家科学基金会资助 PHY~2209424 的支持。

\clearpage
\appendix

\begin{figure}[!th]
    \centering
    \includegraphics[width=\linewidth]{images/O0IBncls_W3_large-nu-test-MEs.pdf}
    \caption{$M_\pi \approx 690$ MeV、Wilson3 涂抹下混合重正化的 $P_z = 1.71$ GeV 矩阵元的两个大-$\nu$ 外推。
    两支外推分别用数据 $z \in [6a, 10a]$（绿）与 $[9a,16a]$（红）拟合。
    各图中的范围由竖直虚线标示。 }
    \label{fig:ME_large_nu_comp}
\end{figure}

\begin{figure}[!th]
    \centering
    \includegraphics[width=\linewidth]{images/O0IBncls_W3_large-nu-test-qPDFs.pdf}
    \caption{$M_\pi \approx 690$ MeV、Wilson3 涂抹数据的大-$\nu$ 外推（分别取 $z \in [6a, 10a]$（绿）与 $[9a,16a]$（红））所产生的准 PDF。}
    \label{fig:qPDF_large_nu_comp}
\end{figure}

\begin{figure}[H]
    \centering
    \includegraphics[width=\linewidth]{images/O0IBncls_W3_large-nu-test-PDFs.pdf}
    \caption{$M_\pi \approx 690$ MeV、Wilson3 涂抹数据的大-$\nu$ 外推（分别取 $z \in [6a, 10a]$（绿）与 $[9a,16a]$（红））所产生的光前 PDF。}
    \label{fig:PDF_large_nu_comp}
\end{figure}

\section{外推拟合范围效应}\label{sec:app_large-nu-comp}
在本节中，我们对最干净的 $M_\pi \approx 690$ MeV Wilson3 数据比较两个不同的大-$\nu$ 外推，以展示结果的稳定性。
在图~\ref{fig:ME_large_nu_comp} 中，我们展示了由范围 $z \in [6a, 10a]$（绿）与 $[9a,16a]$（红）内的数据所作的外推。
我们观察到，两支外推彼此以及与 $\nu \approx 6-20$ 的数据都处于远小于 1$\sigma$ 的偏差之内；然而二者的误差却颇为不同。
对更大距离数据的外推给出大得多的误差，而更短距离的数据可以更强地约束外推。
在正文中，我们采用来自 $z \in [7a, 16a]$ 的外推，以利用短距离数据提供的误差缩减，同时对大 $\nu$ 拟合得更好。

接下来，我们计算傅里叶变换，得到每支外推的准 PDF，示于图~\ref{fig:qPDF_large_nu_comp}。
我们看到中心值符合得非常好。
在大部分范围内，较小的 $z \in [6a,10a]$ 的误差包住了 $z \in [9a,16a]$ 准 PDF 的均值。
同样，较大距离的外推具有大得多的误差棒。
尽管 $z \in [6a, 10a]$ 数据对外推的约束更强，两支准 PDF 的误差中仍存在振荡与夹点。

最后，在图~\ref{fig:PDF_large_nu_comp} 中，我们展示由两支外推产生的光前 PDF。
与准 PDF 类似，二者中心值始终处于彼此的 1$\sigma_{z\in[6a,10a]}$ 之内；但误差相当不同。
这些结果表明，尽管大-$\nu$ 外推不同，PDF 的中心值仍保持相当的稳定性，而误差量化则可能需要更深入的探究。
有一些新技术可用于补充简单的 $e^{-m\nu}/\nu^d$ 外推，包括贝叶斯与机器学习技术；其中一些已为赝 PDF 方法论而发展，但可移植到 LaMET~\cite{Chowdhury:2024ymm,Dutrieux:2024rem,Dutrieux:2025jed}。
这应在未来工作中进一步探索。

\section{大-$\nu$ 外推}\label{sec:app_large-nu}
为完整性起见，我们在本节展示其余的大-$\nu$ 外推（图~\ref{fig:O0IB_large-nu-all}）。
正文结果所用的拟合范围为：奇异与轻核子数据的 Wilson3（HYP5）分别取 $[7a,16a]$ 与 $[7a,12a]$（$[6a,10a]$ 与 $[6a,11a]$）。
在上一节附录中，我们已表明低至 $z = 6a$ 起点的拟合范围是有效选择。
我们看到 Wilson3 数据在长距离处的约束远好于 HYP5。
这导致 HYP5 外推比 Wilson3 衰减得更慢，但我们不预期这种长距离行为会剧烈改变可信区域 $x \sim 0.2-0.8$ 内的物理。

\begin{figure*}[h!]
    \centering
    \includegraphics[width=.48\textwidth]{images/O0IBncls_W3_large-nu.pdf}
    \includegraphics[width=.48\textwidth]{images/O0IBncls_H5_large-nu.pdf}\\
    \includegraphics[width=.48\textwidth]{images/O0IBncll_W3_large-nu.pdf}
    \includegraphics[width=.48\textwidth]{images/O0IBncll_H5_large-nu.pdf}
    \caption{$M_\pi \approx 690$（上排）与 $310$（下排）MeV、Wilson3（左列）与 HYP5（右列）涂抹下混合重正化的 $P_z = 1.71$ GeV 矩阵元的大-$\nu$ 外推。
    各外推用每幅图中两条竖直虚线之间的数据拟合。}
    \label{fig:O0IB_large-nu-all}
\end{figure*}

% ---- 参考文献（自英文版 .bbl 原样嵌入，保留英文） ----

\begin{thebibliography}{10}

\bibitem{Good:2024iur}
William Good, Kinza Hasan, and Huey-Wen Lin.
\newblock {Toward the first gluon parton distribution from the LaMET}.
\newblock {\em J. Phys. G}, 52(3):035105, 2025.

\bibitem{Amoroso:2022eow}
S.~Amoroso et~al.
\newblock {Snowmass 2021 Whitepaper: Proton Structure at the Precision
  Frontier}.
\newblock {\em Acta Phys. Polon. B}, 53(12):12--A1, 2022.

\bibitem{Ji:2013dva}
Xiangdong Ji.
\newblock {Parton Physics on a Euclidean Lattice}.
\newblock {\em Phys. Rev. Lett.}, 110:262002, 2013.

\bibitem{Radyushkin:2017cyf}
A.~V. Radyushkin.
\newblock {Quasi-parton distribution functions, momentum distributions, and
  pseudo-parton distribution functions}.
\newblock {\em Phys. Rev. D}, 96(3):034025, 2017.

\bibitem{Ji:2020ect}
Xiangdong Ji, Yu-Sheng Liu, Yizhuang Liu, Jian-Hui Zhang, and Yong Zhao.
\newblock {Large-momentum effective theory}.
\newblock {\em Rev. Mod. Phys.}, 93(3):035005, 2021.

\bibitem{Constantinou:2020hdm}
Martha Constantinou et~al.
\newblock {Parton distributions and lattice-QCD calculations: Toward 3D
  structure}.
\newblock {\em Prog. Part. Nucl. Phys.}, 121:103908, 2021.

\bibitem{Lin:2023kxn}
Huey-Wen Lin.
\newblock {Overview of Lattice Results for Hadron Structure}.
\newblock {\em Few Body Syst.}, 64(3):58, 2023.

\bibitem{Fan:2020cpa}
Zhouyou Fan, Rui Zhang, and Huey-Wen Lin.
\newblock {Nucleon gluon distribution function from 2 + 1 + 1-flavor lattice
  QCD}.
\newblock {\em Int. J. Mod. Phys. A}, 36(13):2150080, 2021.

\bibitem{Fan:2022kcb}
Zhouyou Fan, William Good, and Huey-Wen Lin.
\newblock {Gluon parton distribution of the nucleon from (2+1+1)-flavor lattice
  QCD in the physical-continuum limit}.
\newblock {\em Phys. Rev. D}, 108(1):014508, 2023.

\bibitem{HadStruc:2021wmh}
Tanjib Khan et~al.
\newblock {Unpolarized gluon distribution in the nucleon from lattice quantum
  chromodynamics}.
\newblock {\em Phys. Rev. D}, 104(9):094516, 2021.

\bibitem{Delmar:2023agv}
Joseph Delmar, Constantia Alexandrou, Krzysztof Cichy, Martha Constantinou, and
  Kyriakos Hadjiyiannakou.
\newblock {Gluon PDF of the proton using twisted mass fermions}.
\newblock {\em Phys. Rev. D}, 108(9):094515, 2023.

\bibitem{HadStruc:2022yaw}
Colin Egerer et~al.
\newblock {Toward the determination of the gluon helicity distribution in the
  nucleon from lattice quantum chromodynamics}.
\newblock {\em Phys. Rev. D}, 106(9):094511, 2022.

\bibitem{Good:2023ecp}
William Good, Kinza Hasan, Allison Chevis, and Huey-Wen Lin.
\newblock {Gluon moment and parton distribution function of the pion from
  Nf=2+1+1 lattice QCD}.
\newblock {\em Phys. Rev. D}, 109(11):114509, 2024.

\bibitem{Salas-Chavira:2021wui}
Alejandro Salas-Chavira, Zhouyou Fan, and Huey-Wen Lin.
\newblock {First glimpse into the kaon gluon parton distribution using lattice
  QCD}.
\newblock {\em Phys. Rev. D}, 106(9):094510, 2022.

\bibitem{Karpie:2019eiq}
Joseph Karpie, Kostas Orginos, Alexander Rothkopf, and Savvas Zafeiropoulos.
\newblock {Reconstructing parton distribution functions from Ioffe time data:
  from Bayesian methods to Neural Networks}.
\newblock {\em JHEP}, 04:057, 2019.

\bibitem{Karpie:2021pap}
Joseph Karpie, Kostas Orginos, Anatoly Radyushkin, and Savvas Zafeiropoulos.
\newblock {The continuum and leading twist limits of parton distribution
  functions in lattice QCD}.
\newblock {\em JHEP}, 11:024, 2021.

\bibitem{Collins:1989gx}
John~C. Collins, Davison~E. Soper, and George~F. Sterman.
\newblock {Factorization of Hard Processes in QCD}.
\newblock {\em Adv. Ser. Direct. High Energy Phys.}, 5:1--91, 1989.

\bibitem{Zhang:2018diq}
Jian-Hui Zhang, Xiangdong Ji, Andreas Sch\"afer, Wei Wang, and Shuai Zhao.
\newblock {Accessing Gluon Parton Distributions in Large Momentum Effective
  Theory}.
\newblock {\em Phys. Rev. Lett.}, 122(14):142001, 2019.

\bibitem{Ji:2020brr}
Xiangdong Ji, Yizhuang Liu, Andreas Sch\"afer, Wei Wang, Yi-Bo Yang, Jian-Hui
  Zhang, and Yong Zhao.
\newblock {A Hybrid Renormalization Scheme for Quasi Light-Front Correlations
  in Large-Momentum Effective Theory}.
\newblock {\em Nucl. Phys. B}, 964:115311, 2021.

\bibitem{Radyushkin:2018cvn}
Anatoly Radyushkin.
\newblock {One-loop evolution of parton pseudo-distribution functions on the
  lattice}.
\newblock {\em Phys. Rev. D}, 98(1):014019, 2018.

\bibitem{LatticePartonLPC:2021gpi}
Yi-Kai Huo et~al.
\newblock {Self-renormalization of quasi-light-front correlators on the
  lattice}.
\newblock {\em Nucl. Phys. B}, 969:115443, 2021.

\bibitem{Zhang:2023bxs}
Rui Zhang, Jack Holligan, Xiangdong Ji, and Yushan Su.
\newblock {Leading power accuracy in lattice calculations of parton
  distributions}.
\newblock {\em Phys. Lett. B}, 844:138081, 2023.

\bibitem{Balitsky:2021qsr}
Ian Balitsky, Wayne Morris, and Anatoly Radyushkin.
\newblock {Short-distance structure of unpolarized gluon pseudodistributions}.
\newblock {\em Phys. Rev. D}, 105(1):014008, 2022.

\bibitem{Balitsky:2019krf}
Ian Balitsky, Wayne Morris, and Anatoly Radyushkin.
\newblock {Gluon Pseudo-Distributions at Short Distances: Forward Case}.
\newblock {\em Phys. Lett. B}, 808:135621, 2020.

\bibitem{Yao:2022vtp}
Fei Yao, Yao Ji, and Jian-Hui Zhang.
\newblock {Connecting Euclidean to light-cone correlations: from flavor
  nonsinglet in forward kinematics to flavor singlet in non-forward
  kinematics}.
\newblock {\em JHEP}, 11:021, 2023.

\bibitem{Hou:2019efy}
Tie-Jiun Hou et~al.
\newblock {New CTEQ global analysis of quantum chromodynamics with
  high-precision data from the LHC}.
\newblock {\em Phys. Rev. D}, 103(1):014013, 2021.

\bibitem{MILC:2013znn}
A.~Bazavov et~al.
\newblock {Lattice QCD Ensembles with Four Flavors of Highly Improved Staggered
  Quarks}.
\newblock {\em Phys. Rev. D}, 87(5):054505, 2013.

\bibitem{Follana:2007rc}
E.~Follana, Q.~Mason, C.~Davies, K.~Hornbostel, G.~P. Lepage, J.~Shigemitsu,
  H.~Trottier, and K.~Wong.
\newblock {Highly improved staggered quarks on the lattice, with applications
  to charm physics}.
\newblock {\em Phys. Rev. D}, 75:054502, 2007.

\bibitem{Bali:2016lva}
Gunnar~S. Bali, Bernhard Lang, Bernhard~U. Musch, and Andreas Sch\"afer.
\newblock {Novel quark smearing for hadrons with high momenta in lattice QCD}.
\newblock {\em Phys. Rev. D}, 93(9):094515, 2016.

\bibitem{Hasenfratz:2001hp}
Anna Hasenfratz and Francesco Knechtli.
\newblock {Flavor symmetry and the static potential with hypercubic blocking}.
\newblock {\em Phys. Rev. D}, 64:034504, 2001.

\bibitem{Luscher:2010iy}
Martin L\"uscher.
\newblock {Properties and uses of the Wilson flow in lattice QCD}.
\newblock {\em JHEP}, 08:071, 2010.
\newblock [Erratum: JHEP 03, 092 (2014)].

\bibitem{NieMiera:2025inn}
Alex NieMiera, William Good, and Huey-Wen Lin.
\newblock {Kaon Gluon Parton Distribution and Momentum Fraction from 2+1+1
  Lattice-QCD with High Statistics}.
\newblock 6 2025.

\bibitem{Anderson:2024evk}
Trey Anderson, W.~Melnitchouk, and N.~Sato.
\newblock {Strangeness in the proton from W+charm production and SIDIS data}.
\newblock 12 2024.

\bibitem{Cerutti:2025yji}
Matteo Cerutti, Alberto Accardi, Ishara~P. Fernando, Shujie Li, Joseph~F.
  Owens, and Sanghwa Park.
\newblock {Systematic uncertainties from higher-twist corrections in DIS at
  large x}.
\newblock {\em Phys. Rev. D}, 111(9):094013, 2025.

\bibitem{Gao:2021dbh}
Xiang Gao, Andrew~D. Hanlon, Swagato Mukherjee, Peter Petreczky, Philipp Scior,
  Sergey Syritsyn, and Yong Zhao.
\newblock {Lattice QCD Determination of the Bjorken-x Dependence of Parton
  Distribution Functions at Next-to-Next-to-Leading Order}.
\newblock {\em Phys. Rev. Lett.}, 128(14):142003, 2022.

\bibitem{Fan:2021bcr}
Zhouyou Fan and Huey-Wen Lin.
\newblock {Gluon parton distribution of the pion from lattice QCD}.
\newblock {\em Phys. Lett. B}, 823:136778, 2021.

\bibitem{Chowdhury:2024ymm}
Talal~Ahmed Chowdhury, Taku Izubuchi, Methun Kamruzzaman, Nikhil Karthik,
  Tanjib Khan, Tianbo Liu, Arpon Paul, Jakob Schoenleber, and Raza~Sabbir
  Sufian.
\newblock {Polarized and unpolarized gluon PDFs: Generative machine learning
  applications for lattice QCD matrix elements at short distance and large
  momentum}.
\newblock {\em Phys. Rev. D}, 111(7):074509, 2025.

\bibitem{Dutrieux:2024rem}
Herv\'e Dutrieux, Joseph Karpie, Kostas Orginos, and Savvas Zafeiropoulos.
\newblock {Simple nonparametric reconstruction of parton distributions from
  limited Fourier information}.
\newblock {\em Phys. Rev. D}, 111(3):034515, 2025.

\bibitem{Dutrieux:2025jed}
Herv\'e Dutrieux, Joe Karpie, Christopher~J. Monahan, Kostas Orginos, Anatoly
  Radyushkin, David Richards, and Savvas Zafeiropoulos.
\newblock {Inverse problem in the LaMET framework}.
\newblock 4 2025.

\bibitem{Edwards:2004sx}
Robert~G. Edwards and Balint Joo.
\newblock {The Chroma software system for lattice QCD}.
\newblock {\em Nucl. Phys. B Proc. Suppl.}, 140:832, 2005.

\end{thebibliography}
